\section[Cote 156-1 : le modèle-segment est un biordre]{Cote 156-1 : le modèle-segment est un biordre\protect\verifie{Folder156_1}}\label{sec:156-1}

La cote 156-1, \emph{[Chapitre] I. Vers une géométrie des formes (topologiques) : notes manuscrites (05/06/1986)}, vingt-six pages datées du 5 au 7~juin 1986, propose une axiomatique purement combinatoire des formes de dimension~$1$ : un ensemble $\mathcal L$ de \emph{lieux}, un ensemble $S$ de \emph{segments}, et pour chaque segment $I$ l'ensemble $\widetilde I \subset \mathcal L$ de ses lieux et la paire $\partial I \subset \widetilde I$ de ses extrémités. La troisième mouture, à partir de la page~5, pose les axiomes F1 à F5 (\lot{156-1}{1}{5}, \lot{156-1}{1}{7}) ; la page~10 appelle \emph{modèle-segment} un modèle où $\mathcal L$ est lui-même un segment, la page~11 y définit, à partir d'une extrémité $x_0$, la relation $x \preccurlyeq y$ par $x \in [x_0, y]$ et affirme que c'est un ordre (\lot{156-1}{1}{11}). La page~12 énonce : \q{Cette relation d'ordre a un plus petit élément, $x_0$, et un plus grand élément, $x_1$. Elle permet de reconstituer la structure de modèle de 1-forme}, les segments étant les $[x,y]$, $x<y$, et l'ordre \q{infiniment divisible}, \q{localement filtrante décroissante}, la clause \q{filtrante croissante} étant ajoutée entre les lignes et lue en partie ; elle esquisse la réciproque sur une ligne qui porte plusieurs mots illisibles (\lot{156-1}{1}{12}). La page~13 conclut qu'\q{une structure de modèle de segment équivaut à une structure de \emph{biordre}} (\lot{156-1}{1}{13}). Aucune des deux directions n'est démontrée sur la page.

\noindent\emph{Conventions.} On écrit les axiomes comme la lecture les énonce, en notant $I_{x,z}$ le sous-segment de $I$ d'extrémités $x, z$ donné par F1 : un segment $J$ avec $\widetilde J \subset \widetilde I$ et $\partial J = \{x,z\}$. On pose $I^\circ = \widetilde I \setminus \partial I$, et un lieu est \emph{régulier} s'il est intérieur à un segment.
\begin{itemize}
\item[F1] Si $x \in \partial I$ et $z \in \widetilde I \setminus \{x\}$, il existe un unique $I_{x,z}$.
\item[F2] $I^\circ \neq \emptyset$ ; si $\partial I = \{x,y\}$ et $z \in I^\circ$, alors $\widetilde I_{x,z} \cap \widetilde I_{y,z} = \{z\}$, et tout segment $J$ avec $\widetilde J \subset \widetilde I$ et $z \in \partial J$ vérifie $\widetilde J \subset \widetilde I_{x,z}$ ou $\widetilde J \subset \widetilde I_{y,z}$.
\item[F3] Si $x \in \partial I \cap \partial J$, il existe $u \in \widetilde I \setminus \{x\}$ et $v \in \widetilde J \setminus \{x\}$ avec $I_{x,u} = J_{x,v}$ ou $\widetilde I_{x,u} \cap \widetilde J_{x,v} = \{x\}$ (la forme de la marge).
\item[F$'$3] Si $x \in \partial I$ et $u, v \in I^\circ$, il existe $w \in I^\circ$ avec $\widetilde I_{x,w} \subset \widetilde I_{x,u} \cap \widetilde I_{x,v}$ (la reconstruction de la lecture, d'après la note du bas de la page~6).
\item[F4] Si $z$ est régulier, $\widetilde I' \cap \widetilde I'' = \{z\}$, $\partial I' = \{x,z\}$ et $\partial I'' = \{y,z\}$, il existe un unique $I$ avec $\widetilde I \supset \widetilde I' \cup \widetilde I''$ et $\partial I = \{x,y\}$.
\item[F5] Un lieu régulier est d'ordre~$2$ : il existe deux segments issus de lui qui ne définissent pas la même branche, et tout segment issu de lui définit la même branche que l'un d'eux.
\end{itemize}
Dans un modèle-segment, $\mathcal L = \widetilde I_0$ avec $\partial I_0 = \{x_0, x_1\}$, on pose $a \leq b$ si $a = b$ ou s'il existe un segment $J$ avec $\partial J = \{x_0, b\}$ et $a \in \widetilde J$ ; par l'unicité dans F1, un tel $J$ est $I_{0\,x_0,b}$, et l'on note $[x_0,b] = \widetilde J$. Un ensemble ordonné est \emph{localement filtrant décroissant} si $x < y_1$ et $x < y_2$ entraînent l'existence de $y$ avec $x < y \leq y_1, y_2$, \emph{localement filtrant croissant} dans l'énoncé dual.

\begin{enonce}[Pages 10 à 13, \emph{le modèle-segment est un biordre}\piste{156-1-segment-biorder}]
Soit un modèle satisfaisant F1, F2, F$'$3 et F4, dont l'ensemble des lieux est un segment $I_0$, $\partial I_0 = \{x_0, x_1\}$. La relation $\leq$ est un ordre sur $\mathcal L$, de plus petit élément $x_0$ et de plus grand élément $x_1$, distincts ; il est dense, localement filtrant décroissant et localement filtrant croissant ; les segments sont exactement les intervalles $[a,b] = \{z \mid a \leq z \leq b\}$, $a < b$, avec $\partial = \{a,b\}$, chacun porté par un seul segment. L'ordre défini à partir de $x_1$ est l'ordre opposé.

Réciproquement, soit $\mathcal L$ un ensemble ordonné ayant un plus petit et un plus grand élément distincts, dense et localement filtrant dans les deux sens. Les intervalles $[a,b]$, $a<b$, avec $\partial [a,b] = \{a,b\}$, satisfont F1 à F5 et F$'$3, $\mathcal L$ est un segment, et l'ordre relu à partir du plus petit élément est l'ordre donné.
\end{enonce}

\begin{proof}[Démonstration du sens direct]
\emph{Deux unicités.} Un segment $J$ avec $\partial J = \{x_0, b\}$ est unique : c'est F1 dans $I_0$, puisque $\widetilde J \subset \mathcal L = \widetilde I_0$. Pour la même raison, deux segments qui ont une extrémité dans $\{x_0, x_1\}$ et les mêmes extrémités sont égaux. Si $u, v \notin \{x_0, x_1\}$ et $\partial J = \partial K = \{u,v\}$, alors $u \in I_0^\circ$ ; F2 dans $I_0$ en $u$ place $\widetilde J$, et de même $\widetilde K$, dans $[x_0, u]$ ou dans $\widetilde I_{0\,x_1,u}$. S'ils tombent dans des sous-segments différents, $v$ est dans leur intersection, qui est $\{u\}$, ce qui est exclu ; s'ils tombent dans le même, F1 dans ce sous-segment, en son extrémité $u$, donne $J = K$. Ainsi \emph{un segment est déterminé par ses extrémités} (le corollaire~3~a de la page~5).

\emph{C'est un ordre.} Si $b \in [x_0, c]$, $b \neq x_0$, F1 dans le segment $[x_0,c]$ fournit un sous-segment d'extrémités $x_0, b$, qui est $[x_0,b]$ ; donc $[x_0, b] \subset [x_0, c]$, d'où la transitivité. Pour l'antisymétrie, soient $a \neq b$ avec $a \in [x_0,b]$ et $b \in [x_0,a]$ ; alors $a \neq x_0$ (aucun segment n'a $\{x_0\}$ pour bord), $[x_0,a] = [x_0,b]$ par ce qui précède, et $a$ est intérieur à $[x_0,b]$. F2 dans $[x_0,b]$ en $a$ dit que $[x_0,a]$ et le sous-segment d'extrémités $b, a$ ne se rencontrent qu'en $a$ ; or $b$ est dans les deux. On a $x_0 \leq a$ puisque $x_0 \in [x_0,a]$, et $a \leq x_1$ puisque $[x_0,x_1] = \mathcal L$ ; $x_0 \neq x_1$ parce que $\partial I_0$ a deux éléments.

\emph{Recollement.} Soient $a$ un lieu régulier, $[x_0,a]$, et $Q$ un segment d'extrémités $a, z$ tel que $[x_0, a] \cap \widetilde Q = \{a\}$. F4 fournit un segment $G$ d'extrémités $x_0, z$ contenant $[x_0,a]$ ; donc $a \in [x_0, z]$, c'est-à-dire $a \leq z$.

\emph{Les extrémités d'un segment sont comparables.} Soit $\partial J = \{u,v\}$. Si l'une est $x_0$ ou $x_1$, c'est clair. Sinon $u \in I_0^\circ$ et F2 place $\widetilde J$ dans $[x_0,u]$, d'où $v \leq u$, ou dans $P = \widetilde I_{0\,x_1,u}$. Dans le second cas, le sous-segment $Q$ de $P$ d'extrémités $u, v$ ne rencontre $[x_0,u]$ qu'en $u$, puisque $[x_0,u] \cap P = \{u\}$, et le recollement donne $u \leq v$.

\emph{Un segment d'extrémités $a < b$ est l'intervalle $[a,b]$.} Si $a = x_0$, c'est la définition. Sinon, soit $K$ ce segment ; il est le sous-segment de $[x_0,b]$ d'extrémités $b, a$, par l'unicité. Le lieu $a$ est intérieur à $[x_0,b]$, et régulier (il est intérieur à $I_0$, car $a \neq x_1$). F2 dans $[x_0,b]$ en $a$ donne $[x_0,a] \cap \widetilde K = \{a\}$, et place tout sous-segment de $[x_0,b]$ issu de $a$ dans $[x_0,a]$ ou dans $\widetilde K$. Si $z \in \widetilde K$, $z \neq a$, le sous-segment de $K$ d'extrémités $a, z$ ne rencontre $[x_0, a]$ qu'en $a$, et le recollement donne $a \leq z$ ; et $z \leq b$ puisque $\widetilde K \subset [x_0,b]$. Inversement, soit $a < z < b$ ; alors $[x_0,z] \subset [x_0,b]$ contient $a$, et le sous-segment de $[x_0,z]$ d'extrémités $z, a$ est un sous-segment de $[x_0,b]$ issu de $a$. S'il était dans $[x_0,a]$, on aurait $z \leq a$, donc $z = a$ ; il est donc dans $\widetilde K$, et $z \in \widetilde K$. Un tel segment existe toujours : F1 dans $[x_0,b]$ en $b$ fournit le sous-segment d'extrémités $b, a$. Avec l'unicité, chaque intervalle $[a,b]$, $a<b$, est porté par un seul segment, et tout segment est un tel intervalle.

\emph{Densité, filtrations.} Si $a < b$, F2 donne un lieu intérieur au segment $[a,b]$, strictement entre $a$ et $b$. Si $a < y_1$ et $a < y_2$, et si aucun des deux n'est $x_1$ (sinon l'autre convient), $y_1$ et $y_2$ sont intérieurs au segment $[a,x_1]$ ; F$'$3 en son extrémité $a$ fournit $w$ intérieur avec $[a,w] \subset [a,y_1] \cap [a,y_2]$, donc $a < w \leq y_1, y_2$. Dualement, F$'$3 appliqué au segment $[x_0,a]$ en son extrémité $a$ donne la filtration croissante.

\emph{Biordre.} Pour l'ordre défini à partir de $x_1$, $a \leq_{1} b$ signifie $a = b$, ou $a$ dans le segment d'extrémités $x_1, b$, qui est l'intervalle $[b, x_1]$ de l'ordre défini à partir de $x_0$ : c'est $b \leq a$.
\end{proof}

\begin{proof}[Démonstration de la réciproque]
Un segment est un couple $a < b$, de lieux $[a,b]$ et de bord $\{a,b\}$ ; il est déterminé par son bord, la condition $a<b$ fixant l'ordre de la paire. F1 : le sous-segment de $[a,b]$ issu de $a$ d'autre extrémité $z$ est $[a,z]$, celui issu de $b$ est $[z,b]$. F2 : $]a,b[$ est non vide par densité ; si $a < z < b$, $[a,z] \cap [z,b] = \{z\}$, et un intervalle contenu dans $[a,b]$ qui a $z$ pour extrémité est $[z,c] \subset [z,b]$ ou $[c,z] \subset [a,z]$. F3 : si deux segments partent de $x$ vers le haut, $[x,b]$ et $[x,d]$, la filtration décroissante donne $y$ avec $x < y \leq b, d$, et $[x,y]$ est un début commun ; vers le bas, de même par la filtration croissante ; si l'un monte et l'autre descend, $[x,b] \cap [c,x] = \{x\}$. F$'$3 : si $x$ est l'extrémité inférieure de $[x,b]$ et $u, v \in\, ]x,b[$, la filtration décroissante donne $w$ avec $x < w \leq u, v$, d'où $[x,w] \subset [x,u] \cap [x,v]$ ; dualement en l'extrémité supérieure. F4 : si $[x,z]$ et $[y,z]$ sont tous deux sous $z$, la filtration croissante donne $w$ avec $x, y \leq w < z$, qui est dans l'intersection, exclu ; de même s'ils sont tous deux au-dessus ; sinon $x < z < y$ (ou l'inverse) et $[x,y]$ est l'unique segment d'extrémités $x,y$, qui contient $[x,z] \cup [z,y]$. F5 : si $a < x < b$, les segments $[a,x]$ et $[x,b]$ n'ont aucun point intérieur commun, donc ne définissent pas la même branche ; un segment $[x,d]$ définit la même branche que $[x,b]$ : la filtration donne $y$ avec $x < y \leq d, b$, la densité $z$ avec $x < z < y$, et $[x,z]$ est le sous-segment commun ; dualement pour $[c,x]$ et $[a,x]$. Enfin $[\bot,\top] = \mathcal L$, et $a \in [\bot,b]$ équivaut à $a \leq b$.
\end{proof}

\begin{remarque}
Le sens direct n'utilise ni F3 ni F5 : il suffit de F1, F2, F$'$3 et F4. F$'$3 est exactement ce qui donne les deux filtrations b) et b$'$), comme le dit l'esquisse de la lecture ; F4, en un lieu régulier, est ce qui fait de chaque intervalle un segment, et il intervient dès la comparabilité des extrémités d'un segment. Dans la réciproque, F4 vaut sans son hypothèse de régularité, et la densité et les deux filtrations suffisent à F1--F5 et F$'$3 ; les bornes ne servent qu'à faire de $\mathcal L$ lui-même un segment.
\end{remarque}

\begin{remarque}
L'ordre n'est pas total en général, comme la lecture le souligne. Soit, sur $\mathbb Q^2$, l'ordre du cône ouvert, $p \prec q$ si $|q_2 - p_2| < q_1 - p_1$, et soit $T$ l'ensemble des points compris, pour cet ordre, entre $(0,0)$ et $(2,0)$. C'est un ensemble ordonné borné, dense (le milieu de $p \prec q$ est strictement entre eux) et localement filtrant dans les deux sens : si $x \prec y_1, y_2$, le point $x + (\varepsilon, 0)$ convient pour $\varepsilon > 0$ assez petit, et dualement $x - (\varepsilon, 0)$. Les points $(1, \tfrac12)$ et $(1, -\tfrac12)$ y sont incomparables. Ses intervalles forment donc un modèle-segment dont l'ordre n'est pas total.
\end{remarque}

\noindent Ne sont pas démontrés ici : les corollaires de F1 et F2 hors du modèle-segment (pages~5 et~6) et la transitivité de la relation \q{même branche} sous F$'$3 ; les corollaires~1 et~2 de F4 (page~7), dont la lecture relève la lacune, et l'équivalence de F5 et F$'$5 (page~8) ; la restriction à une partie (pages~8 et~9) ; l'orientation des modèles réguliers connexes (pages~14 à~17) ; le contre-exemple de la page~21 aux axiomes F1--F5 et l'axiome F6 (page~22) ; les réseaux et leur cas réduit (pages~23 et~24) ; l'étude, inachevée, de deux segments de mêmes extrémités (pages~24 à~26).
